\chapter{Происхождение общей корреляции среды}
\label{ch:v8-baryon-common-environment-correlation-origin-gate}

Предыдущая глава обнаружила отрезок трёхчастичных ковариаций
$-1/2\leq c\leq1$. Проверим, выбирает ли уже построенный микроскопический
родитель полного шума коллективную точку $c=1$.

\section{Второй момент многокопийной среды}

Пусть $F_a^{(r)}$ обозначает скачок канала $a$ на копии $r=1,2,3$, а
$|\eta_{r,a}\rangle$ --- соответствующее возбуждение среды. Вакуумный
второй момент общего взаимодействия равен
\begin{equation}
 \langle0|G^2|0\rangle
 =\sum_a\sum_{r,s=1}^{3}R_{rs}F_a^{(r)}F_a^{(s)},
 \qquad
 R_{rs}=\langle\eta_{r,a}|\eta_{s,a}\rangle.
 \label{eq:v8-baryon-environment-second-moment}
\end{equation}
Однокопийный родитель фиксирует только $R_{rr}=1$. Перестановочная
симметрия сводит оставшуюся свободу к
\begin{equation}
 R=R_c=(1-c)I_3+c\mathbf1\mathbf1^{\mathsf T},
 \qquad -\frac12\leq c\leq1,
 \label{eq:v8-baryon-environment-covariance-family}
\end{equation}
но не выбирает точку отрезка.

Для самосопряжённых скачков генератор можно записать как
\begin{equation}
 \mathcal L_c(X)=-\frac12\sum_{r,s}R_{rs}
 [F^{(r)},[F^{(s)},X]].
 \label{eq:v8-baryon-copy-double-commutator}
\end{equation}
Если $X=X_1\otimes I\otimes I$, все члены, кроме $r=s=1$, исчезают.
Следовательно, ограничение $\mathcal L_c$ на одночастичные наблюдаемые не
зависит от $c$. На двухчастичном наблюдаемом зависимость уже ненулевая;
для точного двухуровневого контроля её квадрат нормы Гильберта--Шмидта
равен $8c^2$.

\section{Два равноправных микроскопических родителя}

Тензорное произведение трёх независимых ячеек среды даёт
\begin{equation}
 R_{
m ind}=I_3,
 \qquad c=0.
 \label{eq:v8-baryon-independent-environment}
\end{equation}
Одна общая ячейка, с которой взаимодействуют все три копии, даёт
\begin{equation}
 R_{\rm com}=\mathbf1\mathbf1^{\mathsf T},
 \qquad c=1.
 \label{eq:v8-baryon-common-environment}
\end{equation}
Оба родителя самосопряжённы, перестановочно-ковариантны и после ограничения
на одну копию возвращают тот же микроскопический контакт Тома VIII.
Калибровочная ковариантность также сохраняется, если метка канала $a$
преобразуется одинаково во всех трёх копиях.

\begin{theorem}[Запрет вывода общей среды из однокопийного родителя]
Однокопийный микроскопический родитель полного шумового генератора не
определяет межкопийные скалярные произведения возбуждений среды. Поэтому он
не различает независимый родитель $c=0$, общий родитель $c=1$ и весь
промежуточный вполне положительный отрезок.
\end{theorem}

\section{Условные способы выбора}

Требование строгой тензорной композиции однокопийных каналов условно
выбирает $c=0$ и нулевое барионное расщепление данного происхождения.
Напротив, отдельная аксиома единой общей среды условно выбирает $c=1$ и
максимальный дискриминатор предыдущей главы. Ни одно из этих требований не
следует из текущего родительского действия.

\begin{nogocriterion}[Остановка барионного спринта]
Коллективная корреляция $c=1$ не выведена. Дальнейшее уточнение численного
барионного остатка без межкопийного двухточечного ядра среды означало бы
выбор точки ковариационного отрезка вручную. Ветка открывается повторно
только после построения общего родительского действия среды либо вывода её
пространственного двухточечного ядра на трёх кварковых копиях.
\end{nogocriterion}

\begin{protocol}
\begin{enumerate}
 \item поле вычислений: \textbf{$\mathbb Q(c)$};
 \item вакуумный первый момент: \textbf{нулевой};
 \item диагональ ковариации: \textbf{единичная};
 \item независимая среда: \textbf{$c=0$, ранг $3$};
 \item общая среда: \textbf{$c=1$, ранг $1$};
 \item граница противокорреляции: \textbf{$c=-1/2$, ранг $2$};
 \item одночастичное ограничение зависит от $c$: \textbf{нет};
 \item двухчастичные наблюдаемые зависят от $c$: \textbf{да};
 \item текущий родитель выбирает $c$: \textbf{нет};
 \item коллективное барионное расщепление: \textbf{не выведено};
 \item статус ветви: \textbf{остановка на точной неединственности}.
\end{enumerate}
\end{protocol}

Машинный сертификат сохранён в
\path{s2t_v8_baryon_common_environment_correlation_origin_gate_results.json}.